\documentclass[10pt,a4paper]{amsart}
\usepackage[margin=19mm]{geometry}
\usepackage{amsmath,amssymb,mathtools}
\usepackage[scr=rsfs,cal=euler]{mathalfa}
\usepackage{xfrac}
\usepackage[unicode,hidelinks,bookmarks=false]{hyperref}
\newcommand{\ie}{i.e.\ }
\newcommand{\cf}{cf.\ }
\newcommand{\resp}{respectively\ }
\newcommand{\Subsection}{\subsection}
\providecommand{\nobreakdash}{\nobreak\hbox{-}}
\setlength{\emergencystretch}{2em}
\multlinegap0pt
\hypersetup{colorlinks,
	linkcolor={black},
	filecolor={black},
	citecolor={blue},
	urlcolor={blue}}  
\newtheorem{theorem}{Theorem}[section]
\newtheorem{remark}{Remark}
\newtheorem{conjecture}{Conjecture}
\newtheorem{Def}{Definition}[section]
\newtheorem{corollary}{Corollary}[section]
\newtheorem{lma}{Lemma}[section]
\numberwithin{equation}{section}
\begin{document}
\title[Arithmetic properties of the 2-color overpartition function ]{Arithmetic properties of the 2-color overpartition function $\overline{p}_\ell(n)$}
\author{H. S. Sumanth Bharadwaj, N. Sujatha; S. Chandankumar}
\date{}
\maketitle
\noindent\emph{Source and licence.} Arithmetic properties of the 2-color overpartition function $\overline{p}_\ell(n)$. Version arXiv:2607.16608v1 (CC BY 4.0). This material is distributed under the Creative Commons Attribution 4.0 International licence, http://creativecommons.org/licenses/by/4.0/, as recorded in the arXiv OAI licence metadata for this exact version and shown on the abstract page https://arxiv.org/abs/2607.16608v1. Full rights notice: licenses/arxiv-2607.16608-CC-BY-4.0.txt.\par\smallskip
\noindent\textit{Editorial scope.} Counted unit: the original \texttt{theorem} environment \texttt{t1} at source lines 238--250 together with its complete proof at lines 251--266; the delivered document keeps source lines 79--292 so that every referenced label and every citation resolves inside it. Native editable source is the paper's own concatenated TeX; the AMS wrapper is editorial formatting only. No publisher class, upstream script or unrelated artwork is redistributed.
	\begin{equation}
		\sum_{n=0}^{\infty} \overline{p}_\ell(n)q^n
		=
		\frac{(-q;q)_\infty(-q^\ell;q^\ell)_\infty}
		{(q;q)_\infty(q^\ell;q^\ell)_\infty}=
		\frac{f_2f_{2\ell}}{f_1^2f_\ell^2}.
		\label{pt}
	\end{equation}
	For example, $\overline{p}_3(3)=10$; the ten such $2$-color overpartitions of $3$ are
	\[
	3_a,\;
	\overline{3}_a,\;
	3_b,\;
	\overline{3}_b,\;
	2_a+1_a,\;
	\overline{2}_a+1_a,\;
	2_a+\overline{1}_a,\;
	\overline{2}_a+\overline{1}_a,\;
	1_a+1_a+1_a,\;
	\overline{1}_a+1_a+1_a.
	\]
	Arithmetic congruences modulo small powers of $2$ and $3$ for $\overline{p}_3(n)$ were investigated by M. S. Mahadeva Naika, S. Shivaprasada Nayaka, and C. Shivashankar in \cite{MSMSSNCS}. The general function $\overline{p}_\ell(n)$, and in particular the way its congruences depend on $\ell$, has not otherwise been treated.
	
	The primary objective of this paper is to give such a treatment, systematically investigating the arithmetic properties of $\overline{p}_{\ell}(n)$ for $\ell=2k,3k,4,6,8,$ and $9$. Our results are of four kinds. First, we prove \emph{$\ell$-uniform families}: a single dissection of the generating function \eqref{pt} yields, for \emph{every} $\ell\ge1$, the congruences of Theorem~\ref{t1} and Corollaries~\ref{c1}--\ref{c2}, so that one identity settles infinitely many functions at once. Second, we obtain \emph{internal congruences} relating a function to a dilate of itself, such as $\overline{p}_{3\ell}(3n+1)\equiv\overline{p}_{3\ell}\!\left(4^\alpha(3n+1)\right)\pmod 4$. Third, for the individual values $\ell=4,6,8,9$ we establish prime-power congruences; the strongest of these is the modulus-$512$ family
	\[
	\overline{p}_{4}(32n+28)\equiv0\pmod{512}\qquad(n\ge0),
	\]
	together with a family of related congruences modulo $128$ and $256$. Fourth, combining the support of $f_1$ and $f_1^3$ on pentagonal and triangular numbers with a quadratic-residue argument modulo an auxiliary prime $p$, we obtain two new vanishing congruences,
	\[
	\overline{p}_{3\ell}(24pn+24b+1)\equiv0\pmod4,\qquad
	\overline{p}_{4}(16pn+16b+2)\equiv0\pmod8,
	\]
	valid for \emph{every} prime $p\ge5$ and every $b$ for which $24b+1$, respectively $8b+1$, is a quadratic non-residue modulo $p$; since infinitely many such $(p,b)$ exist, this yields infinitely many further congruences at a single stroke.
	
	The methods employed throughout the paper are elementary and rely chiefly on classical $q$-series identities, dissections, and generating function manipulations; no appeal is made to the theory of modular forms. Section 2 collects the dissection lemmas used throughout. Section 3 proves the $\ell$-uniform families and their corollaries, together with the internal congruences of Theorem~\ref{t1}. Section 4 treats the individual cases $\ell=4,6,8,9$ in turn, including the two quadratic-residue vanishing congruences described above. Section 5 records further observations and conjectures.
	
	\section{Preliminaries}
	\noindent We begin this section by introducing Ramanujan's general theta-function $f(a,b)$, defined as:
	\begin{align}
		f(a,b): = \sum_{n=-\infty}^{\infty} a^{\frac{n(n+1)}{2}} b^{\frac{n(n-1)}{2}}, \quad |ab| < 1.
	\end{align}
	The following definitions of theta-functions $\varphi$, $\psi$ and $f$ are classical:
	\begin{eqnarray}
		\varphi(q)&:=&f(q,q)= \frac{f_2^5}{f_1^2f_4^2},\\
		\psi(q)&:=&f(q,q^3) =\frac{f_2^2}{f_1},\\
		f(-q)&:=&f(-q,-q^2)=f_1.
	\end{eqnarray}
	\begin{lma}
		Suppose $l \geq 1$, $m \geq 1$ are any integer, and $p$ is any prime. Then we have
		\begin{align}\label{bt}
			f^{p^{l-1}}_{pm} \equiv f_{m}^{p^{l}} \pmod {p^l}.
		\end{align}
	\end{lma}
	\begin{lma}\cite{NBK}\label{l2}
		The following $2$-dissection hold
		\begin{align}
			\dfrac{1}{f_1f_3}&=\dfrac{f_8^2f_{12}^5}{f_2^2f_4f_6^4f_{24}^2}+q\dfrac{f_4^5f_{24}^2}{f_2^4f_6^2f_8^2f_{12}}.\label{1byf1f3}
		\end{align}
	\end{lma}
	\begin{lma}[{\cite{xia2013analogues}}]\label{l2b}
		The following $2$-dissection holds modulo $3$:
		\begin{align}
			\frac{f_1^2\,f_6^4}{f_2^2\,f_3^2}
			\equiv
			\frac{f_4^2\,f_{12}^4}{f_2\,f_6\,f_8\,f_{24}}
			+q\,\frac{f_8\,f_{12}\,f_{24}}{f_4}
			\pmod{3}.\label{f12byf33}
		\end{align}
	\end{lma}
	\begin{lma} [{\cite[Entry~25]{berndt2012ramanujan}}] The following 2-dissections hold: 
		\begin{equation}\label{f12}
			f_1^2 =\frac{f_2 f_8^5}{f_4^2 f_{16}^2}
			-2q\frac{f_2 f_{16}^2}{f_8},
		\end{equation}
		\begin{equation}\label{f14}
			f_1^4 =\frac{f_4^{10}}{f_2^2 f_{8}^4}
			-4q\frac{f_2^2 f_{8}^4}{f_4^2},
		\end{equation}
		\begin{equation}\label{1byf12}
			\frac{1}{f_1^2} =\frac{f_8^5}{f_2^5 f_{16}^2}
			+2q\frac{f_4^2 f_{16}^2}{f_2^5 f_8},
		\end{equation}
		\begin{equation}\label{1byf14}
			\frac{1}{f_1^4} =\frac{f_4^{14}}{f_2^{14} f_{8}^4}
			+4q\frac{f_4^2 f_{8}^4}{f_2^{10}}.
		\end{equation}
	\end{lma}
	\begin{lma}[{\cite{hirschhorn1993cubic}}]The following 2-dissection hold:
		\begin{align}
			\frac{f_3^3}{f_1} &= \frac{f_4^3f_6^2}{f_2^2f_{12}}+q\frac{f_{12}^3}{f_4}.\label{f33byf1}
			%\frac{f_3}{f_1^3} &= \frac{f_4^6f_6^3}%{f_2^9f_{12}^2}+3q\frac{f_4^2f_6f_{12}^2}{f_2^7}.\label{f3byf13}
		\end{align}
	\end{lma}
	
	\begin{lma}[{\cite{mdb}}] The following 3-dissection hold:
		\begin{align}
			\frac{f_4}{f_1}=\frac{f_{12}f_{18}^{4}}{f_3^{3}f_{36}^{2}}+q\,\frac{f_6^{2}f_9^{3}f_{36}}{f_3^{4}f_{18}^{2}}+2q^{2}\,\frac{f_6f_{18}f_{36}}{f_3^{3}}.\label{f4byf1}
	\end{align}\end{lma}
	\begin{lma}[{\cite{mdb}}] The following 3-dissections hold:
		\begin{align}
			&\frac{f_2}{f_1^2} = \frac{f_6^4 f_9^6}{f_3^8 f_{18}^3} + 2q \frac{f_6^3 f_9^3}{f_3^7} + 4q^2 \frac{f_6^2 f_{18}^3}{f_3^6},\label{f2byf12}\\
			&\frac{f_1^2}{f_2} = \frac{f_9^2}{f_{18}} - 2q \frac{f_3 f_{18}^2}{f_6 f_9},\label{f12byf2}\\
			&\frac{f_2^2}{f_1} = \frac{f_6f_9^2}{f_3f_{18}} +q \frac{ f_{18}^2}{ f_9}.\label{f22byf1}
	\end{align}\end{lma}
	\begin{lma} [{\cite{hirschhorn2014congruence}}] The following 3-dissection hold:
		
		\begin{align}
			f_1f_2=\frac{f_6\,f_9^4}{f_3f_{18}^2}
			-qf_9f_{18}-2q^2\frac{f_3\,f_{18}^4}{f_6\,f_9^2}.\label{f1f2}
		\end{align}
		%	\begin{align}
			%		\frac{1}{f_1f_2} &= \frac{f_9^9}{f_3^6 f_6^2 f_{18}^3} + q \frac{f_9^6}{f_3^5 f_6^3} + 3q^2 - 2q^3 \frac{f_{18}^6}{f_3^3 f_6^5} + 4q^4 \frac{f_{18}^9}{f_3^2 f_6^6f_9^3}.\label{1byf1f2}
			%		\end{align}
	\end{lma}
	
	\begin{lma} [{\cite{xia2013analogues}}] The following 2-dissection hold:
		\begin{align}      
			\frac{f_3^4}{f_1^4}
			&=\frac{f_4^8 f_6^2 f_{12}^{4}}{f_2^{10} f_8^2 f_{24}^{2}}
			+ 4q \frac{f_4^5 f_6^3 f_{12}}{f_2^9}
			+ 4q^2 \frac{f_4^2 f_6^4 f_8^2 f_{24}^{2}}{f_2^8 f_{12}^{2}}.\label{f34f14}
		\end{align}
	\end{lma}
	\begin{lma}[{\cite[Theorem ~2.2]{cui2013}}]\label{a1}
		For any prime $p\ge5$,
		\begin{equation}\label{a17}
			f_1=\sum\limits_{\substack{k=\frac{1-p}{2}\\k\neq\frac{\pm p-1}{6}
			}}^{\frac{p-1}{2}}{(-1)^kq^{\frac{3k^2+k}{2}}f\left(-q^{\frac{3p^2+(6k+1)p}{2}},-q^{\frac{3p^2-(6k+1)p}{2}}\right)}+(-1)^{\frac{\pm p-1}{6}}q^{\frac{p^2-1}{24}}f_{p^2},
		\end{equation}
		where
		\begin{equation}
			\dfrac{\pm p-1}{6}:=\begin{cases}
				\frac{p-1}{6}, & \text{if $p \equiv 1 \pmod{6}$},\\
				\frac{-p-1}{6}, & \text{if $p \equiv -1 \pmod{6}$}.\nonumber
			\end{cases}
		\end{equation}
	\end{lma}
	
	\begin{lma}[{\cite[Theorem ~2.1]{cui2013}}]
		For any odd prime p,
		\begin{align}
			\psi(q)=\sum_{m=0}^{\frac{p-3}{2}}q^{\frac{m^{2}+m}{2}}f\left ( q^\frac{p^2+(2m+1)p}{2},q^\frac{p^2-(2m+1)p}{2} \right )+q^{\frac{p^{2}-1}{8}}\frac{f_{2p^{2}}^{2}}{f_{p^{2}}}.\label{a2}
		\end{align}
	\end{lma}
	Also,
	\begin{equation*}
		\frac{m^{2}+m}{2}\not \equiv \frac{p^{2}-1}{8} \pmod{p} {\hspace{1mm}} \text{for} {\hspace{1mm}} 0\leq m\leq \frac{p-3}{2}.
	\end{equation*}
	
	\begin{lma}[{\cite[Theorem 2.1]{LW}}]
		For any odd prime $p$,
		\begin{equation}
			f_{1}^{3}=\sum_{k=-\frac{p-1}{2}}^{\frac{p-3}{2}}(-1)^kq^{\frac{k^{2}+k}{2}}B_k(q^p)+(-1)^{\frac{p-1}{2}}pq^{\frac{p^{2}-1}{8}}f_{p^{2}}^{3},\label{abc1}
		\end{equation}
	\end{lma}
	where $$B_k(q):=\frac{1}{2}\sum_{n=-\infty}^{\infty}(-1)^n(2pn+2k+1)q^{\dfrac{pn^2+(2k+1)n}{2}}.$$
	
	\section{Congruences for $\overline{p}_{2\ell}(n)$ and $\overline{p}_{3\ell}(n)$}
	In the following theorem, we derive the 2-dissections and 3-dissections of the function $\overline{p}_{2\ell}(n)$ and $\overline{p}_{3\ell}(n)$ respectively. These dissections play a pivotal role in establishing the arithmetic properties of the partition function $\overline{p}_{\ell}(n)$.

	\begin{theorem}\label{t1} For all $n\geq0$ and $\ell\geq1$,
		\begin{align}
			\sum_{n=0}^{\infty}\overline{p}_{2\ell}(2n)q^n=\frac{f_{2\ell}f_4^5}{f_{\ell}^2f_1^4f_{8}^2}=\frac{\varphi(q^2)}{\varphi(-q^\ell)\varphi^2(-q)},\label{res4}\\
			\sum_{n=0}^{\infty}\overline{p}_{2\ell}(2n+1)q^n=2\,\frac{f_{2\ell}f_2^2f_{8}^2}
			{f_{\ell}^2f_1^4f_4} =\frac{\psi(q^4)}{\varphi(-q^\ell)\varphi^2(-q)},\label{res5}\\
			\sum_{n=0}^{\infty}\overline{p}_{3\ell}(3n)q^n=
			\frac{f_{2\ell}f_2^4f_3^6}{f_{\ell}^2f_1^8f_{6}^3}=\frac{\varphi^3(-q^3)}{\varphi(-q^\ell)\varphi^4(-q)},\label{res1}\\
			\sum_{n=0}^{\infty}\overline{p}_{3\ell}(3n+1)q^n=
			2\,\frac{f_{2\ell}f_2^3f_3^3}{f_{\ell}^2f_1^7},\label{res2}\\
			\sum_{n=0}^{\infty}\overline{p}_{3\ell}(3n+2)q^n=
			4\,\frac{f_{2\ell}f_2^2f_{6}^3}{f_{\ell}^2f_1^6}.\label{res3}
		\end{align}
	\end{theorem}

	\begin{proof}
		Invoking \eqref{f2byf12} in \eqref{pt} with $\ell=3\ell$, we have
		\begin{align}\label{r1}
			\sum_{n=0}^{\infty}\overline{p}_{3\ell}(n)q^n=
			\frac{f_{6\ell}f_6^4f_9^6}{f_{3\ell}^2f_3^8f_{18}^3}
			+2q\,\frac{f_{6\ell}f_6^3f_9^3}{f_{3\ell}^2f_3^7}
			+4q^2\,\frac{f_{6\ell}f_6^2f_{18}^3}{f_{3\ell}^2f_3^6}.
		\end{align}
		Equations \eqref{res1}--\eqref{res3} directly follows by extracting the terms involving $q^{3n}$, $q^{3n+1}$ and $q^{3n+2}$ respectively from the equation \eqref{r1}. Now, invoking \eqref{1byf12} in \eqref{pt} with $\ell=2\ell$, we have
		\begin{align}
			\sum_{n=0}^{\infty}\overline{p}_{2\ell}(n)q^n=\frac{f_{4\ell}f_8^5}{f_{2\ell}^2f_2^4f_{16}^2}
			+2q\,\frac{f_{4\ell}f_4^2f_{16}^2}
			{f_{2\ell}^2f_2^4f_8}. \label{r2}
		\end{align}
		Equations \eqref{res4} and \eqref{res5} are obtained  by extracting the terms involving $q^{2n}$ and $q^{2n+1}$ respectively from the equation \eqref{r2}. This completes the proof.
	\end{proof}

	%The following Corollary \ref{c1} and Corollary \ref{c2}, follows directly from the above Theorem \ref{t1}
	\begin{corollary}\label{c1} Suppose $t$ and $\ell$ are natural numbers and $4|t$ and $9|\ell$, then
		\begin{align}
			\overline{p}_{2t}(4n+2)\equiv 0\pmod{4},\label{cor1}\\
			\overline{p}_{2t}(4n+3)\equiv 0\pmod{8},\label{cor2}\\
			\overline{p}_{3\ell}(9n+3j)\equiv 0\pmod{8}, \,\,\text{for}\,\,j=1,2.\label{cor3}
			%\overline{p}_{3t}(6n+5)\equiv 0\pmod{8}.\label{cor4}
		\end{align}
	\end{corollary}
	\begin{proof}
		These congruences follow on specializing $t$ and $\ell$ in \eqref{res4}, \eqref{res5}, and \eqref{res1}, reducing modulo the stated modulus by \eqref{bt}, and extracting the relevant arithmetic progression.
	\end{proof}
	\begin{corollary}\label{c2} For all $n\geq0$ and $\ell\geq1$,
		\begin{align}
			\sum_{n=0}^{\infty}\overline{p}_{2\ell}(2n)q^n\equiv \frac{1}{\varphi(-q^2)\varphi(-q^\ell)}\pmod{4},\label{res6}\\
			\sum_{n=0}^{\infty}\overline{p}_{3\ell}(3n)q^n\equiv
			\frac{\varphi^3(-q^3)}{\varphi(-q^\ell)}\pmod{8},\label{res8}\\
			\sum_{n=0}^{\infty}\overline{p}_{2\ell}(2n+1)q^n\equiv2\,\frac{\psi(q^4)}
			{\varphi(-q^\ell)}\pmod{8}, \label{res7}\\
			\overline{p}_{3\ell}(3n+1)\equiv 0 \pmod{2},\\
			\overline{p}_{3\ell}(3n+2)\equiv0 \pmod{4}.
		\end{align}
	\end{corollary}
	\begin{proof}
		Each congruence follows on reducing the dissections of Theorem~\ref{t1} modulo the stated modulus by \eqref{bt}.
	\end{proof}

\begin{thebibliography}{99}
\bibitem[LW]{LW} Wang, L.: Arithmetic properties of $7$-regular partitions, {\em Ramanujan J.} \textbf{47} (2018), 99--115.
\bibitem[MSMSSNCS]{MSMSSNCS} Mahadeva Naika, M. S., Nayaka, S. S. and Shivashankar, C.: Infinite families of congruences for $2$-color overpartitions, {\em Tamsui Oxford J. Inf. Math. Sci.} \textbf{30} (2014), 61--79.
\bibitem[NBK]{NBK} Baruah, N. D. and Ojah, K. K.: Partitions with designated summands in which all parts are odd, {\em Integers} \textbf{15} (2015), Article A53.
\bibitem[berndt2012ramanujan]{berndt2012ramanujan} Berndt, B. C.: {\em Ramanujan's Notebooks, Part III}, Springer, New York, 2012.
\bibitem[cui2013]{cui2013} Cui, S.-P. and Gu, N. S. S.: Arithmetic properties of $l$-regular partitions, {\em Adv. Appl. Math.} \textbf{51} (2013), 507--523.
\bibitem[hirschhorn1993cubic]{hirschhorn1993cubic} Hirschhorn, M. D., Garvan, F. G. and Borwein, J. M.: Cubic analogs of the Jacobi cubic theta function $\theta(z,q)$, {\em Canad. J. Math.} \textbf{45} (1993), 673--694.
\bibitem[hirschhorn2014congruence]{hirschhorn2014congruence} Hirschhorn, M. D. and Sellers, J. A.: A congruence modulo $3$ for partitions into distinct non-multiples of four, {\em J. Integer Seq.} \textbf{17} (2014), Article 14.9.5.
\bibitem[mdb]{mdb} Hirschhorn, M. D.: {\em The Power of $q$}, Developments in Mathematics, Vol. 49, Springer, Cham, 2017.
\bibitem[xia2013analogues]{xia2013analogues} Xia, E. X. and Yao, O. M.: Analogues of Ramanujan's partition identities, {\em Ramanujan J.} \textbf{31} (2013), 373--396.
\end{thebibliography}
\end{document}
